\section{Aspectos éticos y protección de datos}
\label{chap:gestion_etica}

Los registros clínicos utilizados en este trabajo fueron cedidos por el Hospital Universitario Central de Asturias (HUCA) previamente anonimizados, con el fin de preservar la privacidad de los pacientes. La protección de estos datos es especialmente relevante en aplicaciones de inteligencia artificial en el ámbito sanitario. En este sentido, el uso de Gemma~4, un modelo de pesos abiertos, permite realizar el procesamiento y la inferencia en un entorno local, sin necesidad de transmitir información clínica a servicios externos.

Desde el punto de vista clínico, la aplicación planteada no pretende sustituir el diagnóstico médico, sino contribuir a priorizar los ECG potencialmente sospechosos para su revisión por especialistas. En este contexto, los falsos negativos podrían retrasar la identificación de pacientes que requieren una evaluación más detallada, mientras que los falsos positivos podrían incrementar innecesariamente la carga asistencial. Por ello, cualquier aplicación práctica requeriría una validación clínica rigurosa, prestando especial atención a la sensibilidad del sistema y manteniendo siempre la supervisión y responsabilidad del personal sanitario \cite{zeppenfeld2022,wilde2023scbt}.
